\section{Orders of vanishing}
\label{section-orders-of-vanishing}

\begin{lemma}
\label{lemma-ord-additive}
Let $R$ be a semi-local Noetherian ring of dimension $1$.
If $a, b \in R$ are nonzerodivisors then
$$
\text{length}_R(R/(ab)) =
\text{length}_R(R/(a)) +
\text{length}_R(R/(b))
$$
and these lengths are finite.
\end{lemma}

\begin{proof}
We saw the finiteness in Lemma \ref{lemma-finite-length-global}.
Additivity holds since there is a short exact sequence
$0 \to R/(a) \to R/(ab) \to R/(b) \to 0$ where the first map
is given by multiplication by $b$. (Use length is additive,
see Lemma \ref{lemma-length-additive}.)
\end{proof}

\begin{definition}
\label{definition-ord}
Suppose that $K$ is a field, and $R \subset K$ is a
local\footnote{We could also define this when $R$ is only
semi-local but this is probably never really what you want!}
Noetherian subring of dimension $1$ with fraction field $K$.
In this case we define the {\it order of vanishing along $R$}
$$
\text{ord}_R : K^* \longrightarrow \mathbf{Z}
$$
by the rule
$$
\text{ord}_R(x) = \text{length}_R(R/(x))
$$
if $x \in R$ and we set
$\text{ord}_R(x/y) = \text{ord}_R(x) - \text{ord}_R(y)$
for $x, y \in R$ both nonzero.
\end{definition}

\noindent
We can use the order of vanishing to compare lattices in a
vector space. Here is the definition.

\begin{definition}
\label{definition-lattice}
Let $R$ be a Noetherian local domain of dimension $1$ with
fraction field $K$. Let $V$ be a finite dimensional $K$-vector space.
A {\it lattice in $V$} is a finite $R$-submodule $M \subset V$ such
that $V = K \otimes_R M$.
\end{definition}

\noindent
The condition $V = K \otimes_R M$ signifies that $M$ contains a
basis for the vector space $V$. We remark that in many places in the
literature the notion of a lattice may be defined only in case the
ring $R$ is a discrete valuation ring. If $R$ is a discrete valuation
ring then any lattice is a free $R$-module, and this may not be the case
in general.

\begin{lemma}
\label{lemma-compare-lattices}
Let $R$ be a Noetherian local domain of dimension $1$ with
fraction field $K$. Let $V$ be a finite dimensional $K$-vector space.
\begin{enumerate}
\item If $M$ is a lattice in $V$ and $M \subset M' \subset V$
is an $R$-submodule of $V$ containing $M$
then the following are equivalent
\begin{enumerate}
\item $M'$ is a lattice,
\item $\text{length}_R(M'/M)$ is finite, and
\item $M'$ is finitely generated.
\end{enumerate}
\item If $M$ is a lattice in $V$ and $M' \subset M$ is an $R$-submodule
of $M$ then $M'$ is a lattice if and only if
$\text{length}_R(M/M')$ is finite.
\item If $M$, $M'$ are lattices in $V$, then so are
$M \cap M'$ and $M + M'$.
\item If $M \subset M' \subset M'' \subset V$ are lattices in $V$
then
$$
\text{length}_R(M''/M) =
\text{length}_R(M'/M) +
\text{length}_R(M''/M').
$$
\item If $M$, $M'$, $N$, $N'$ are lattices in $V$ and
$N \subset M \cap M'$, $M + M' \subset N'$, then we have
\begin{eqnarray*}
& & \text{length}_R(M/M \cap M') - \text{length}_R(M'/M \cap M')\\
& = &
\text{length}_R(M/N) - \text{length}_R(M'/N) \\
& = &
\text{length}_R(M + M' / M') - \text{length}_R(M + M'/M) \\
& = &
\text{length}_R(N' / M') - \text{length}_R(N'/M)
\end{eqnarray*}
\end{enumerate}
\end{lemma}

\begin{proof}
Proof of (1). Assume (1)(a). Say $y_1, \ldots, y_m$ generate $M'$.
Then each $y_i = x_i/f_i$ for some $x_i \in M$ and
nonzero $f_i \in R$.
Hence we see that $f_1 \ldots f_m M' \subset M$.
Since $R$ is Noetherian local of dimension $1$
we see that $\mathfrak m^n \subset (f_1 \ldots f_m)$
for some $n$ (for example combine
Lemmas \ref{lemma-one-equation} and
Proposition \ref{proposition-dimension-zero-ring} or combine
Lemmas \ref{lemma-finite-length} and \ref{lemma-length-infinite}).
In other words $\mathfrak m^nM' \subset M$ for some $n$
Hence
$\text{length}(M'/M) < \infty$ by Lemma \ref{lemma-length-finite},
in other words (1)(b) holds.
Assume (1)(b). Then $M'/M$ is a finite $R$-module
(see Lemma \ref{lemma-finite-length-finite}).
Hence $M'$ is a finite $R$-module as an extension of finite $R$-modules.
Hence (1)(c). The implication
(1)(c) $\Rightarrow$ (1)(a) follows from the remark following
Definition \ref{definition-lattice}.

\medskip\noindent
Proof of (2). Suppose
$M$ is a lattice in $V$ and $M' \subset M$ is an $R$-submodule.
We have seen in (1) that if $M'$ is a lattice, then
$\text{length}_R(M/M') < \infty$. Conversely, assume that
$\text{length}_R(M/M') < \infty$. Then $M'$ is finitely generated
as $R$ is Noetherian and for some $n$ we have
$\mathfrak m^n M \subset M'$ (Lemma \ref{lemma-length-infinite}).
Hence it follows
that $M'$ contains a basis for $V$, and $M'$ is a lattice.

\medskip\noindent
Proof of (3). Assume $M$, $M'$ are lattices in $V$.
Since $R$ is Noetherian the submodule $M \cap M'$ of $M$ is finite.
As $M$ is a lattice we can find
$x_1, \ldots, x_n \in M$ which form a $K$-basis for
$V$. Because $M'$ is a lattice we can write $x_i = y_i/f_i$ with
$y_i \in M'$ and $f_i \in R$. Hence $f_ix_i \in M \cap M'$. Hence
$M \cap M'$ is a lattice also.
The fact that $M + M'$ is a lattice follows from part (1).

\medskip\noindent
Part (4) follows from additivity of lengths
(Lemma \ref{lemma-length-additive})
and the exact sequence
$$
0 \to M'/M \to M''/M \to M''/M' \to 0
$$
Part (5) follows from repeatedly applying part (4).
\end{proof}

\begin{definition}
\label{definition-distance}
Let $R$ be a Noetherian local domain of dimension $1$ with
fraction field $K$. Let $V$ be a finite dimensional $K$-vector space.
Let $M$, $M'$ be two lattices in $V$. The {\it distance between
$M$ and $M'$} is the integer
$$
d(M, M') = \text{length}_R(M/M \cap M') - \text{length}_R(M'/M \cap M')
$$
of Lemma \ref{lemma-compare-lattices} part (5).
\end{definition}

\noindent
In particular, if $M' \subset M$, then
$d(M, M') = \text{length}_R(M/M')$.

\begin{lemma}
\label{lemma-properties-distance-function}
Let $R$ be a Noetherian local domain of dimension $1$ with
fraction field $K$. Let $V$ be a finite dimensional $K$-vector space.
This distance function has the property that
$$
d(M, M'') = d(M, M') + d(M', M'')
$$
whenever given three lattices $M$, $M'$, $M''$ of $V$.
In particular we have $d(M, M') = - d(M', M)$.
\end{lemma}

\begin{proof}
Omitted.
\end{proof}

\begin{lemma}
\label{lemma-order-vanishing-determinant}
Let $R$ be a Noetherian local domain of dimension $1$ with
fraction field $K$. Let $V$ be a finite dimensional $K$-vector space.
Let $\varphi : V \to V$ be a $K$-linear isomorphism.
For any lattice $M \subset V$ we have
$$
d(M, \varphi(M)) = \text{ord}_R(\det(\varphi))
$$
\end{lemma}

\begin{proof}
We can see that the integer $d(M, \varphi(M))$ does not depend
on the lattice $M$ as follows. Suppose that $M'$ is a second such
lattice. Then we see that
\begin{eqnarray*}
d(M, \varphi(M)) & = & d(M, M') + d(M', \varphi(M)) \\
& = & d(M, M') + d(\varphi(M'), \varphi(M)) + d(M', \varphi(M'))
\end{eqnarray*}
Since $\varphi$ is an isomorphism we see that
$d(\varphi(M'), \varphi(M)) = d(M', M) = -d(M, M')$, and hence
$d(M, \varphi(M)) = d(M', \varphi(M'))$. Moreover, both sides of the
equation (of the lemma) are additive in $\varphi$, i.e.,
$$
\text{ord}_R(\det(\varphi \circ \psi))
=
\text{ord}_R(\det(\varphi))
+
\text{ord}_R(\det(\psi))
$$
and also
\begin{eqnarray*}
d(M, \varphi(\psi((M))) & = &
d(M, \psi(M)) + d(\psi(M), \varphi(\psi(M))) \\
& = & d(M, \psi(M)) + d(M, \varphi(M))
\end{eqnarray*}
by the independence shown above. Hence it suffices to prove the lemma
for generators of $\text{GL}(V)$. Choose an isomorphism
$K^{\oplus n} \cong V$. Then $\text{GL}(V) = \text{GL}_n(K)$ is
generated by elementary matrices $E$.
The result is clear for $E$ equal to the identity matrix.
If $E = E_{ij}(\lambda)$ with $i \not = j$, $\lambda \in K$,
$\lambda \not = 0$, for example
$$
E_{12}(\lambda) =
\left(
\begin{matrix}
1 & \lambda & \ldots \\
0 & 1 & \ldots \\
\ldots & \ldots & \ldots
\end{matrix}
\right)
$$
then with respect to a different basis we get $E_{12}(1)$.
The result is clear for $E = E_{12}(1)$ by taking as lattice
$R^{\oplus n} \subset K^{\oplus n}$. Finally, if $E = E_i(a)$,
with $a \in K^*$ for example
$$
E_1(a) =
\left(
\begin{matrix}
a & 0 & \ldots \\
0 & 1 & \ldots \\
\ldots & \ldots & \ldots
\end{matrix}
\right)
$$
then $E_1(a)(R^{\oplus b}) = aR \oplus R^{\oplus n - 1}$ and
it is clear that $d(R^{\oplus n}, aR \oplus R^{\oplus n - 1})
= \text{ord}_R(a)$ as desired.
\end{proof}

\begin{lemma}
\label{lemma-finite-extension-dim-1}
Let $A \to B$ be a ring map. Assume
\begin{enumerate}
\item $A$ is a Noetherian local domain of dimension $1$,
\item $A \subset B$ is a finite extension of domains.
\end{enumerate}
Let $L/K$ be the corresponding finite extension of fraction fields.
Let $y \in L^*$ and $x = \text{Nm}_{L/K}(y)$.
In this situation $B$ is semi-local.
Let $\mathfrak m_i$, $i = 1, \ldots, n$ be the maximal ideals of $B$.
Then
$$
\text{ord}_A(x) =
\sum\nolimits_i
[\kappa(\mathfrak m_i) : \kappa(\mathfrak m_A)]
\text{ord}_{B_{\mathfrak m_i}}(y)
$$
where $\text{ord}$ is defined as in Definition \ref{definition-ord}.
\end{lemma}

\begin{proof}
The ring $B$ is semi-local by Lemma \ref{lemma-finite-in-codim-1}.
Write $y = b/b'$ for some $b, b' \in B$.
By the additivity of $\text{ord}$ and multiplicativity of
$\text{Nm}$ it suffices to prove the lemma for
$y = b$ or $y = b'$. In other words we may assume $y \in B$.
In this case the right hand side of the formula is
$$
\sum [\kappa(\mathfrak m_i) : \kappa(\mathfrak m_A)]
\text{length}_{B_{\mathfrak m_i}}((B/yB)_{\mathfrak m_i})
$$
By Lemma \ref{lemma-pushdown-module} this is equal to
$\text{length}_A(B/yB)$. By Lemma \ref{lemma-order-vanishing-determinant}
we have
$$
\text{length}_A(B/yB) = d(B, yB) =
\text{ord}_A(\det\nolimits_K(L \xrightarrow{y} L)).
$$
Since $x = \text{Nm}_{L/K}(y) = \det\nolimits_K(L \xrightarrow{y} L)$
by definition the lemma is proved.
\end{proof}














